% 来源及取材范围见分题文件；此为整理者转录的正文，不是下载原件。
\begin{definition}[Proper map]
Suppose $f\colon X\to Y$ is a continuous map between two topological spaces. Then $f$ is a \emph{proper map} if for every compact $K\subset\subset Y$, the set $f^{-1}(K)$ is compact.
\end{definition}
\begin{theorem}[Rothstein 1935]\label{thm:Rothstein}
There exists no proper holomorphic mapping of the unit bidisc $\D^2=\D\times\D\subset\C^2$ to the unit ball $\bB_2\subset\C^2$.
\end{theorem}
\begin{definition}[Analytic disc]
A nonconstant holomorphic mapping $\varphi\colon\D\to\C^n$ is called an \emph{analytic disc}. If the mapping $\varphi$ extends continuously to the closed unit disc $\widebar\D$, then the mapping $\varphi\colon\widebar\D\to\C^n$ is called a \emph{closed analytic disc}.
Often we call the image $\Delta=\varphi(\D)$ the analytic disc rather than the mapping. For a closed analytic disc we write $\partial\Delta=\varphi(\partial\D)$ and call it the boundary of the analytic disc.
\end{definition}
\begin{proposition}\label{prop:noanaldiscinsphere}
The unit sphere $S^{2n-1}=\partial\bB_n\subset\C^n$ contains no analytic discs.
\end{proposition}
\begin{proof}
Suppose there is a holomorphic function $g\colon\D\to\C^n$ such that the image $g(\D)$ is contained in the unit sphere. In other words, for all $z\in\D$,
\[\snorm{g(z)}^2=\sabs{g_1(z)}^2+\sabs{g_2(z)}^2+\cdots+\sabs{g_n(z)}^2=1.\]
Without loss of generality (after composing with a unitary matrix), assume that $g(0)=(1,0,0,\ldots,0)$. Consider the first component and notice that $g_1(0)=1$. If a sum of nonnegative numbers is less than or equal to $1$, then they all are. Hence, $\sabs{g_1(z)}\leq1$. The maximum principle says that $g_1(z)=1$ for all $z\in\D$. But then $g_k(z)=0$ for all $k=2,\ldots,n$ and all $z\in\D$. Therefore, $g$ is constant and thus not an analytic disc.
\end{proof}
Before we prove the theorem, let us make the statement about proper maps taking boundary to boundary precise.
\begin{lemma}\label{lemma:bndrytobndry}
Let $U\subset\R^n$ and $V\subset\R^m$ be bounded domains and let $f\colon U\to V$ be continuous. Then $f$ is proper if and only if for every sequence $\{p_k\}$ in $U$ such that $p_k\to p\in\partial U$, the set of limit points of $\bigl\{f(p_k)\bigr\}$ lies in $\partial V$.
\end{lemma}
\begin{proof}
Suppose $f$ is proper. Let $\{p_k\}$ be a sequence in $U$ such that $p_k\to p\in\partial U$. Take any convergent subsequence $\bigl\{f(p_{k_\ell})\bigr\}$ of $\bigl\{f(p_k)\bigr\}$ converging to some $q\in\widebar V$. Consider $E=\bigl\{f(p_{k_\ell})\bigr\}$ as a set. Let $\widebar E$ be the closure of $E$ in $V$ (subspace topology). The inverse image $f^{-1}(\widebar E)$ is not compact (it contains a sequence going to $p\in\partial U$) and hence $\widebar E=(E\cup\{q\})\cap V$ is not compact either as $f$ is proper. That means that $q\notin V$ and so $q\in\partial V$. As we took an arbitrary convergent subsequence of $\bigl\{f(p_k)\bigr\}$, $q$ was an arbitrary limit point. Therefore, all limit points are in $\partial V$.

Let us prove the converse. Suppose that for every sequence $\{p_k\}$ in $U$ such that $p_k\to p\in\partial U$, the set of limit points of $\bigl\{f(p_k)\bigr\}$ lies in $\partial V$. Take a closed set $E\subset V$ (subspace topology) and suppose $f^{-1}(E)$ is not compact. Then there exists a sequence $\{p_k\}$ in $f^{-1}(E)$ such that $p_k\to p\in\partial U$, because $f^{-1}(E)$ is closed (in $U$), bounded, but not compact. The hypothesis then says that the limit points of $\bigl\{f(p_k)\bigr\}$ are in $\partial V$. Hence, $E$ has limit points in $\partial V$ and is not compact.
\end{proof}
\begin{proof}[Proof of Rothstein's theorem]
Suppose there is a proper holomorphic map $f\colon\D^2\to\bB_2$. Fix some $e^{i\theta}$ in the boundary of the disc $\D$. Take a sequence $w_k\in\D$ such that $w_k\to e^{i\theta}$. The functions $g_k(\zeta)=f(\zeta,w_k)$ map the unit disc into $\bB_2$. By Montel's theorem and by passing to a subsequence, assume that the sequence of functions converges (uniformly on compact subsets) to a limit $g\colon\D\to\overline{\bB}_2$. As $(\zeta,w_k)\to(\zeta,e^{i\theta})\in\partial\D^2$, by Lemma~\ref{lemma:bndrytobndry}, $g(\D)\subset\partial\bB_2$, and hence $g$ is constant by Proposition~\ref{prop:noanaldiscinsphere}.

Let $g_k'$ denote the derivative (we differentiate each component). The functions $g_k'$ converge to $g'=0$. So for an arbitrary fixed $\zeta\in\D$, $\frac{\partial f}{\partial z_1}(\zeta,w_k)\to0$. This limit holds for all $e^{i\theta}$ and some subsequence of an arbitrary sequence $\{w_k\}$ where $w_k\to e^{i\theta}$. The holomorphic mapping $w\mapsto\frac{\partial f}{\partial z_1}(\zeta,w)$, therefore, extends continuously to the closure $\widebar\D$ and is zero on $\partial\D$. We apply the maximum principle or the Cauchy formula and the fact that $\zeta$ was arbitrary to find $\frac{\partial f}{\partial z_1}\equiv0$. By symmetry $\frac{\partial f}{\partial z_2}\equiv0$. Therefore, $f$ is constant, which is a contradiction as $f$ was proper.
\end{proof}
